\section{Evaluation}
\label{sec:eval}

The evaluation reported here consists of two verification case studies. Its
purpose is to establish that the compile-time structure \sysname{} recovers is
sufficient to move real prefill work off the critical path, and to locate exactly
which call sites benefit and which do not. It is not a general performance study;
Section~\ref{sec:threats} states the scope precisely.

\subsection{Setup}

Both case studies run Qwen3-4B~\cite{yang2025qwen3} on a single NVIDIA RTX 3090.
\sysname{} is built on vLLM~\cite{kwon2023pagedattention}, and the baseline is
therefore the same byLLM runtime served by vLLM with its automatic prefix caching
(APC) enabled. We compare three configurations:

\begin{itemize}\itemsep2pt
\item \textbf{Baseline}: the same byLLM runtime with vLLM's automatic prefix
      caching.
\item \textbf{Oracle}: everything cached, no prefill required. Oracle is
      \emph{measured}, not modelled: an identically-sized prompt is replayed
      twice and the second serve is read, yielding a 100\% cache hit.
\item \textbf{\sysname{}}: compiler-derived speculative prefill, with no probe
      decode and no deploy-time prewarm.
\end{itemize}

Prefill--decode aggregation is disabled in these runs because the 3090 cannot
sustain it for this model. All reported values are median TTFT: for case study 1,
over 20 cold-start requests; for case study 2, over the calls that 20 multi-hop
questions issue at each call site. The measurements are first requests, not
steady-state requests behind a warm cache.

\subsection{Case study 1: multi-intent supervisor with four subagents}

The first topology, shown in Figure~\ref{fig:topology}, is a multi-intent agent
with a supervisor dominating four subagents. An entry stage normalizes the
request, a routing stage selects one of four candidate desks through an
LLM-driven graph visit, and the selected desk runs a three-stage pipeline whose
middle stages are separated by retrieval. Each desk terminates in a typed brief.
The retrieval steps and the tool calls between stages are what open the idle
windows \sysname{} spends. We run 20 problems sampled from public datasets
including GSM8K~\cite{cobbe2021gsm8k}.

\begin{figure}[t]
\centering
% Case-study topology: supervisor with four subagents.
{\setlength{\fboxsep}{0pt}\setlength{\unitlength}{1mm}
\begin{picture}(84,50)(0,0)

  % supervisor
  \obox{0}{22}{15}{6}{normalize}
  \put(15,25){\vector(1,0){5}}
  \obox{20}{22}{15}{6}{route}
  \put(20,18){\makebox(15,4)[l]{\scriptsize\itshape visit by llm}}
  \put(20,14.5){\makebox(15,4)[l]{\scriptsize\itshape select 1 of 4}}

  % fan-out spine
  \put(35,25){\line(1,0){3}}
  \put(38,4){\line(0,1){42}}

  % MathDesk
  \put(38,46){\vector(1,0){3}}
  \obox{41}{43}{11}{6}{analyze}
  \put(52,46){\vector(1,0){3}}
  \obox{55}{43}{14}{6}{solve}
  \put(69,46){\vector(1,0){3}}
  \obox{72}{43}{12}{6}{finalize}
  \put(41,39.5){\makebox(43,3.5)[l]{\scriptsize MathDesk\quad$\rightarrow$ MathBrief}}

  % CodeDesk
  \put(38,32){\vector(1,0){3}}
  \obox{41}{29}{11}{6}{analyze}
  \put(52,32){\vector(1,0){3}}
  \obox{55}{29}{14}{6}{implement}
  \put(69,32){\vector(1,0){3}}
  \obox{72}{29}{12}{6}{finalize}
  \put(41,25.5){\makebox(43,3.5)[l]{\scriptsize CodeDesk\quad$\rightarrow$ CodeBrief}}

  % DocDesk
  \put(38,18){\vector(1,0){3}}
  \obox{41}{15}{11}{6}{analyze}
  \put(52,18){\vector(1,0){3}}
  \obox{55}{15}{14}{6}{answer}
  \put(69,18){\vector(1,0){3}}
  \obox{72}{15}{12}{6}{finalize}
  \put(41,11.5){\makebox(43,3.5)[l]{\scriptsize DocDesk\quad$\rightarrow$ DocBrief}}

  % WriteDesk
  \put(38,4){\vector(1,0){3}}
  \obox{41}{1}{11}{6}{analyze}
  \put(52,4){\vector(1,0){3}}
  \obox{55}{1}{14}{6}{draft}
  \put(69,4){\vector(1,0){3}}
  \obox{72}{1}{12}{6}{finalize}
  \put(41,-2.5){\makebox(43,3.5)[l]{\scriptsize WriteDesk\quad$\rightarrow$ WriteBrief}}

\end{picture}}

\vspace{2mm}
\caption{Case study 1: a supervisor routes to one of four subagents, each a
three-stage pipeline separated by retrieval.}
\label{fig:topology}
\end{figure}

Table~\ref{tab:percomponent} reports per-component TTFT, and
Figure~\ref{fig:percomponent} shows the same data graphically. The pattern is
sharp and it is the pattern the design predicts.

\begin{table*}[t]
\centering
\caption{Case study 1: per-component median TTFT over 20 cold-start requests,
Qwen3-4B on one RTX 3090. Oracle is measured by replaying one identically-sized
prompt twice and reading the second serve. ``Headroom captured'' is the fraction
of the baseline-to-Oracle gap that \sysname{} closes.}
\label{tab:percomponent}
\footnotesize
\setlength{\tabcolsep}{4pt}
\begin{tabular}{@{}lrrrrrrl@{}}
\toprule
component & $n$ & prompt tok & baseline & \sysname{} & Oracle & vs.\ baseline & \multicolumn{1}{r@{}}{headroom capt.} \\
\midrule
\texttt{normalize (entry)}    & 20 &  193 &  37.3\,ms & \textbf{37.6\,ms} & 35.3\,ms & 0.99$\times$ & \multicolumn{1}{r@{}}{n/a (at floor)} \\
\texttt{route (visit-by-llm)} & 20 &  549 &  65.8\,ms & \textbf{65.9\,ms} & 37.8\,ms & 1.00$\times$ & \multicolumn{1}{r@{}}{0\%} \\
\texttt{MathDesk.analyze}     &  8 &  500 &  65.8\,ms & \textbf{66.1\,ms} & 37.8\,ms & 1.00$\times$ & \multicolumn{1}{r@{}}{0\%} \\
\texttt{CodeDesk.analyze}     &  4 &  423 &  64.0\,ms & \textbf{64.3\,ms} & 34.5\,ms & 1.00$\times$ & \multicolumn{1}{r@{}}{0\%} \\
\texttt{DocDesk.analyze}      &  3 &  283 &  52.3\,ms & \textbf{52.3\,ms} & 35.2\,ms & 1.00$\times$ & \multicolumn{1}{r@{}}{0\%} \\
\texttt{WriteDesk.analyze}    &  5 &  294 &  52.3\,ms & \textbf{52.5\,ms} & 36.6\,ms & 1.00$\times$ & \multicolumn{1}{r@{}}{0\%} \\
\texttt{MathDesk.solve}       &  8 &  879 & 119.8\,ms & \textbf{67.8\,ms} & 37.7\,ms & 1.77$\times$ & \multicolumn{1}{r@{}}{63\%} \\
\texttt{CodeDesk.implement}   &  4 &  835 & 117.9\,ms & \textbf{49.7\,ms} & 36.2\,ms & 2.37$\times$ & \multicolumn{1}{r@{}}{83\%} \\
\texttt{DocDesk.answer}       &  3 &  707 & 102.3\,ms & \textbf{62.3\,ms} & 37.1\,ms & 1.64$\times$ & \multicolumn{1}{r@{}}{61\%} \\
\texttt{WriteDesk.draft}      &  5 &  676 & 102.0\,ms & \textbf{50.7\,ms} & 37.9\,ms & 2.01$\times$ & \multicolumn{1}{r@{}}{80\%} \\
\texttt{MathDesk.finalize}    &  8 & 1453 & 216.3\,ms & \textbf{42.7\,ms} & 40.1\,ms & 5.07$\times$ & \multicolumn{1}{r@{}}{99\%} \\
\texttt{CodeDesk.finalize}    &  4 & 1292 & 187.0\,ms & \textbf{42.6\,ms} & 37.2\,ms & 4.39$\times$ & \multicolumn{1}{r@{}}{96\%} \\
\texttt{DocDesk.finalize}     &  3 & 1296 & 185.0\,ms & \textbf{41.5\,ms} & 39.2\,ms & 4.46$\times$ & \multicolumn{1}{r@{}}{98\%} \\
\texttt{WriteDesk.finalize}   &  5 & 1258 & 171.1\,ms & \textbf{41.9\,ms} & 39.6\,ms & 4.09$\times$ & \multicolumn{1}{r@{}}{98\%} \\
\midrule
\textbf{sum per request}      & 20 & $\sim$3300 & \textbf{469.9\,ms} & \textbf{265.1\,ms} & \textbf{186.2\,ms} & \textbf{1.77$\times$} & \multicolumn{1}{r@{}}{\textbf{72\%}} \\
\bottomrule
\end{tabular}
\end{table*}

\begin{figure}[t]
\centering
% generated by gen_figures.py from story.md
{\footnotesize\setlength{\tabcolsep}{2pt}
\begin{tabular}{@{}l@{\hspace{3pt}}l@{\hspace{2pt}}l@{}}
\texttt{\scriptsize normalize} & {\color{cbase}\rule{4.48mm}{1.5mm}} & \scriptsize 37.3 \\[0.35mm]
 & {\color{cours}\rule{4.52mm}{1.5mm}} & \scriptsize 37.6 \\[0.35mm]
 & {\color{coracle}\rule{4.24mm}{1.5mm}} & \scriptsize 35.3 \\[1.4mm]
\texttt{\scriptsize route} & {\color{cbase}\rule{7.91mm}{1.5mm}} & \scriptsize 65.8 \\[0.35mm]
 & {\color{cours}\rule{7.92mm}{1.5mm}} & \scriptsize 65.9 \\[0.35mm]
 & {\color{coracle}\rule{4.54mm}{1.5mm}} & \scriptsize 37.8 \\[1.4mm]
\texttt{\scriptsize Math.analyze} & {\color{cbase}\rule{7.91mm}{1.5mm}} & \scriptsize 65.8 \\[0.35mm]
 & {\color{cours}\rule{7.95mm}{1.5mm}} & \scriptsize 66.1 \\[0.35mm]
 & {\color{coracle}\rule{4.54mm}{1.5mm}} & \scriptsize 37.8 \\[1.4mm]
\texttt{\scriptsize Code.analyze} & {\color{cbase}\rule{7.69mm}{1.5mm}} & \scriptsize 64.0 \\[0.35mm]
 & {\color{cours}\rule{7.73mm}{1.5mm}} & \scriptsize 64.3 \\[0.35mm]
 & {\color{coracle}\rule{4.15mm}{1.5mm}} & \scriptsize 34.5 \\[1.4mm]
\texttt{\scriptsize Doc.analyze} & {\color{cbase}\rule{6.29mm}{1.5mm}} & \scriptsize 52.3 \\[0.35mm]
 & {\color{cours}\rule{6.29mm}{1.5mm}} & \scriptsize 52.3 \\[0.35mm]
 & {\color{coracle}\rule{4.23mm}{1.5mm}} & \scriptsize 35.2 \\[1.4mm]
\texttt{\scriptsize Write.analyze} & {\color{cbase}\rule{6.29mm}{1.5mm}} & \scriptsize 52.3 \\[0.35mm]
 & {\color{cours}\rule{6.31mm}{1.5mm}} & \scriptsize 52.5 \\[0.35mm]
 & {\color{coracle}\rule{4.40mm}{1.5mm}} & \scriptsize 36.6 \\[1.4mm]
\texttt{\scriptsize Math.solve} & {\color{cbase}\rule{14.40mm}{1.5mm}} & \scriptsize 119.8 \\[0.35mm]
 & {\color{cours}\rule{8.15mm}{1.5mm}} & \scriptsize 67.8 \\[0.35mm]
 & {\color{coracle}\rule{4.53mm}{1.5mm}} & \scriptsize 37.7 \\[1.4mm]
\texttt{\scriptsize Code.implement} & {\color{cbase}\rule{14.17mm}{1.5mm}} & \scriptsize 117.9 \\[0.35mm]
 & {\color{cours}\rule{5.97mm}{1.5mm}} & \scriptsize 49.7 \\[0.35mm]
 & {\color{coracle}\rule{4.35mm}{1.5mm}} & \scriptsize 36.2 \\[1.4mm]
\texttt{\scriptsize Doc.answer} & {\color{cbase}\rule{12.30mm}{1.5mm}} & \scriptsize 102.3 \\[0.35mm]
 & {\color{cours}\rule{7.49mm}{1.5mm}} & \scriptsize 62.3 \\[0.35mm]
 & {\color{coracle}\rule{4.46mm}{1.5mm}} & \scriptsize 37.1 \\[1.4mm]
\texttt{\scriptsize Write.draft} & {\color{cbase}\rule{12.26mm}{1.5mm}} & \scriptsize 102.0 \\[0.35mm]
 & {\color{cours}\rule{6.09mm}{1.5mm}} & \scriptsize 50.7 \\[0.35mm]
 & {\color{coracle}\rule{4.56mm}{1.5mm}} & \scriptsize 37.9 \\[1.4mm]
\texttt{\scriptsize Math.finalize} & {\color{cbase}\rule{26.00mm}{1.5mm}} & \scriptsize 216.3 \\[0.35mm]
 & {\color{cours}\rule{5.13mm}{1.5mm}} & \scriptsize 42.7 \\[0.35mm]
 & {\color{coracle}\rule{4.82mm}{1.5mm}} & \scriptsize 40.1 \\[1.4mm]
\texttt{\scriptsize Code.finalize} & {\color{cbase}\rule{22.48mm}{1.5mm}} & \scriptsize 187.0 \\[0.35mm]
 & {\color{cours}\rule{5.12mm}{1.5mm}} & \scriptsize 42.6 \\[0.35mm]
 & {\color{coracle}\rule{4.47mm}{1.5mm}} & \scriptsize 37.2 \\[1.4mm]
\texttt{\scriptsize Doc.finalize} & {\color{cbase}\rule{22.24mm}{1.5mm}} & \scriptsize 185.0 \\[0.35mm]
 & {\color{cours}\rule{4.99mm}{1.5mm}} & \scriptsize 41.5 \\[0.35mm]
 & {\color{coracle}\rule{4.71mm}{1.5mm}} & \scriptsize 39.2 \\[1.4mm]
\texttt{\scriptsize Write.finalize} & {\color{cbase}\rule{20.57mm}{1.5mm}} & \scriptsize 171.1 \\[0.35mm]
 & {\color{cours}\rule{5.04mm}{1.5mm}} & \scriptsize 41.9 \\[0.35mm]
 & {\color{coracle}\rule{4.76mm}{1.5mm}} & \scriptsize 39.6 \\[1.4mm]
\midrule[0.4pt]
\textbf{\scriptsize sum / request} & {\color{cbase}\rule{56.48mm}{1.5mm}} & \scriptsize 469.9 \\[0.35mm]
 & {\color{cours}\rule{31.87mm}{1.5mm}} & \scriptsize 265.1 \\[0.35mm]
 & {\color{coracle}\rule{22.38mm}{1.5mm}} & \scriptsize 186.2 \\[1.4mm]
\end{tabular}}
\\[1mm]{\scriptsize{\color{cbase}\rule{3mm}{1.5mm}}~baseline (vLLM APC)\quad{\color{cours}\rule{3mm}{1.5mm}}~\sysname\quad{\color{coracle}\rule{3mm}{1.5mm}}~Oracle\quad(median TTFT, ms)}

\caption{Per-component median TTFT for case study 1 (data of
Table~\ref{tab:percomponent}). Gains concentrate in the late stages, whose
prompts are assembled from values produced earlier in the request.}
\label{fig:percomponent}
\end{figure}

\paragraph{Where there is no gain, and why.} The entry stage
(\texttt{normalize}) is already at the floor: at 193 prompt tokens its baseline
TTFT of 37.3\,ms is 2.0\,ms above the 35.3\,ms Oracle, so there is essentially
nothing to remove. The routing stage and all four \texttt{analyze} stages show
1.00$\times$ and capture 0\% of their headroom, even though their Oracle values
show that headroom exists --- \texttt{route}, for instance, sits at 65.8\,ms
against a 37.8\,ms Oracle. This is the expected consequence of running without the
probe mechanism of Section~\ref{sec:design}: until the supervisor has emitted its
routing decision, the four desk-entry prompts are equally reachable and none of
them is worth committing the idle window to.

\paragraph{Where the gain is, and why.} The middle stages
(\texttt{solve}, \texttt{implement}, \texttt{answer}, \texttt{draft}) improve by
1.64--2.37$\times$ and capture 61--83\% of their headroom. Their prompts contain
the output of the preceding \texttt{analyze} stage, which provenance identifies
as an argument of a statically known consumer, and the retrieval step between the
two opens the window in which to prefill it. The terminal stages improve most:
4.09--5.07$\times$, capturing 96--99\% of headroom, landing within
2.3--5.4\,ms of their respective Oracle values. These are the longest prompts in the workflow
(1258--1453 tokens) because they accumulate everything the desk has produced, and
almost all of that content is available before the call is issued.

\paragraph{Cost where there is no benefit.} On the six components that do not
improve, \sysname{} is never more than 0.3\,ms slower than the baseline
(37.3$\rightarrow$37.6, 65.8$\rightarrow$65.9, 65.8$\rightarrow$66.1,
64.0$\rightarrow$64.3, 52.3$\rightarrow$52.3, 52.3$\rightarrow$52.5). Speculative
prefill that does not pay off does not visibly cost, at least at this scale.

\paragraph{Aggregate.} Summed over the components a single request executes,
\(\Sigma\)TTFT falls from 469.9\,ms to 265.1\,ms, a 1.77$\times$ reduction, against
a measured Oracle of 186.2\,ms. \sysname{} therefore captures 72\% of the total
headroom that any KV-reuse system could obtain on this hardware and workload.

\subsection{Case study 2: a ReAct tool loop}

The second topology tests a structure the first does not contain: a
data-dependent loop. It is a single-agent ReAct-style
loop~\cite{yao2023react}, shown in Figure~\ref{fig:react}, that alternates
reasoning with retrieval for two to five turns before emitting a typed brief.

The compiler recovers the loop back-edge --- \texttt{integrate} may be followed
by \texttt{reason}, \texttt{integrate}, or \texttt{finalize} --- and, more
usefully, a self-referential return across it: the running notes are the return
value of \texttt{integrate} and simultaneously an argument of the next
\texttt{reason}, of the next \texttt{integrate}, and of \texttt{finalize}. That
is a provenance fact, not a topology fact, and it is what makes the loop
prefillable at all.

We run 20 multi-hop questions over a closed synthetic corpus of 37 passages
constructed so that no passage contains both ends of a hop, forcing each hop to
issue its own retrieval.

\begin{figure}[t]
\centering
% Case-study 2 topology: single-agent ReAct loop with a data-dependent back-edge.
{\setlength{\fboxsep}{0pt}\setlength{\unitlength}{1mm}
\begin{picture}(84,29)(0,0)

  \obox{2}{12}{14}{6}{reason}
  \put(16,15){\vector(1,0){6}}
  \put(14,18.2){\makebox(10,3.4)[c]{\scriptsize SEARCH}}
  \obox{22}{12}{20}{6}{retrieve (RAG)}
  \put(22,8){\makebox(20,4)[c]{\scriptsize $\approx$0.3\,s gap}}
  \put(42,15){\vector(1,0){6}}
  \obox{48}{12}{16}{6}{integrate}

  % back-edge
  \put(64,15){\line(1,0){4}}
  \put(68,15){\line(0,1){9}}
  \put(68,24){\line(-1,0){59}}
  \put(9,24){\vector(0,-1){6}}
  \put(20,24.6){\makebox(30,4)[c]{\scriptsize loop, 2--5 turns}}

  % exit
  \put(9,12){\line(0,-1){6}}
  \put(9,6){\vector(1,0){43}}
  \put(12,2){\makebox(20,4)[l]{\scriptsize ANSWER}}
  \obox{52}{3}{26}{6}{finalize : QABrief}

\end{picture}}

\caption{Case study 2: a single-agent ReAct loop with a data-dependent back-edge.
The retrieval step opens an idle window of roughly 0.3\,s per turn.}
\label{fig:react}
\end{figure}

\begin{table}[t]
\centering
\caption{Case study 2: median TTFT per call site, Qwen3-4B on one RTX 3090, 20
multi-hop questions. The loop is data-dependent, so the two configurations do not
execute the same number of calls; call counts are given for both.}
\label{tab:react}
\footnotesize
\setlength{\tabcolsep}{3pt}
\begin{tabular}{@{}lrrrrr@{}}
\toprule
call site & calls (b/o) & base & \sysname{} & ratio & cached tok \\
\midrule
\texttt{reason}    & 74 / 84 & 39.3\,ms & 39.3\,ms & 1.00$\times$ & 448\,$\rightarrow$\,448 \\
\texttt{integrate} & 58 / 71 & 39.4\,ms & 38.3\,ms & 1.03$\times$ & 368\,$\rightarrow$\,544 \\
\textbf{\texttt{finalize}} & \textbf{20 / 20} & \textbf{66.8\,ms} & \textbf{39.2\,ms} & \textbf{1.71$\times$} & \textbf{304\,$\rightarrow$\,496} \\
\bottomrule
\end{tabular}
\end{table}

\begin{figure}[t]
\centering
% generated by gen_figures.py from story.md
{\footnotesize\setlength{\tabcolsep}{2pt}
\begin{tabular}{@{}l@{\hspace{3pt}}l@{\hspace{2pt}}l@{}}
\texttt{\scriptsize reason} & {\color{cbase}\rule{15.30mm}{1.5mm}} & \scriptsize 39.3 \\[0.35mm]
 & {\color{cours}\rule{15.30mm}{1.5mm}} & \scriptsize 39.3 \\[1.4mm]
\texttt{\scriptsize integrate} & {\color{cbase}\rule{15.34mm}{1.5mm}} & \scriptsize 39.4 \\[0.35mm]
 & {\color{cours}\rule{14.91mm}{1.5mm}} & \scriptsize 38.3 \\[1.4mm]
\texttt{\scriptsize finalize} & {\color{cbase}\rule{26.00mm}{1.5mm}} & \scriptsize 66.8 \\[0.35mm]
 & {\color{cours}\rule{15.26mm}{1.5mm}} & \scriptsize 39.2 \\[1.4mm]
\end{tabular}}
\\[0.5mm]{\scriptsize{\color{cbase}\rule{3mm}{1.5mm}}~baseline\quad{\color{cours}\rule{3mm}{1.5mm}}~\sysname\quad(median TTFT, ms)}

\caption{Case study 2 median TTFT per call site (data of
Table~\ref{tab:react}).}
\label{fig:reactbars}
\end{figure}

Table~\ref{tab:react} reports the result, and Figure~\ref{fig:reactbars} plots
it. Three observations follow.

First, \texttt{reason} shows no improvement (39.3\,ms in both configurations, with
cached tokens unchanged at 448). There is no idle window immediately preceding
it: control passes from \texttt{integrate} straight into \texttt{reason} with no
tool call in between, so there is no time in which to prefill anything for it.

Second, \texttt{finalize} improves from 66.8\,ms to 39.2\,ms, a 1.71$\times$
reduction, with cached tokens rising from 304 to 496. It consumes the accumulated
notes that \texttt{integrate} produced, and those notes are prefilled during the
retrieval gap, so the terminal call arrives at an already-warm prefix.

Third, \texttt{integrate} improves only slightly, from 39.4\,ms to 38.3\,ms
(1.03$\times$), even though its cached tokens rise substantially, from 368 to 544.
It too consumes a value produced by \texttt{reason} and prefilled in the gap, but
the incremental prompt it adds per turn is small, so the absolute time removed is
correspondingly small.

The loop's data dependence means the two configurations executed different
numbers of \texttt{reason} and \texttt{integrate} calls (74 versus 84, and 58
versus 71); the comparison in Table~\ref{tab:react} is therefore between
per-call medians and not between matched call sequences. The \texttt{finalize}
row, which is the row carrying the gain, is matched at 20 calls in both
configurations.
